\section{The all-width oracle for smooth non-Gaussian priors}
\label{sec:nongaussian-oracle}
Let $\mathcal G_{C,B}$ consist of one-hidden-layer tanh networks
\[
 g(x)=\sum_{j=1}^{C'} c_j\tanh(a_j\cdot(1,x)),
 \qquad C'\le C,\qquad
 |c_j|\le B,\qquad\|a_j\|_\infty\le B.
\]
An output bias can be represented by one additional constant hidden
feature; this changes only the width convention by one. No separate
bound on the teacher output is assumed.

\begin{theorem}[A smooth non-Gaussian shallow-comparison oracle]
\label{thm:nongaussian-shallow-oracle}
Fix the coordinate law $P$ satisfying
\eqref{eq:coordinate-curvature}, the input dimension $d$, the teacher
coefficient cap $B\ge1$, and accuracy/confidence
$\varepsilon,\delta\in(0,1)$. Let $X$ have any distribution on the
unit sphere in $\mathbb R^d$. Labels may be bounded by a fixed
constant, or may have a fixed marginal sub-Gaussian envelope.

For every hidden depth $H\ge1$, width $N\ge C$, and $0<\nu\le1$,
use the independent coordinate prior of Section~\ref{sec:nongaussian-law},
tanh hidden activation, a linear raw readout, and the likelihood
\begin{equation}
 d\Pi_{S_n,\lambda}(w)=Z_{S_n,\lambda}^{-1}
 \exp\!\left[-\frac{N}{2\nu\lambda}
                  \widehat R_{S_n}(z_w)\right]d\pi_{N,H,\nu}(w).
 \label{eq:nongaussian-likelihood}
\end{equation}
There are constants $a,c>0$, depending only on the fixed quantities
above, and a deterministic calibration specified below, such that
\begin{equation}
 n\ge(2+C)^c\exp\!\left\{c[\sqrt H+(1-\nu)H]\right\}
 \label{eq:nongaussian-sample-budget}
\end{equation}
implies, with probability at least $1-\delta$,
\begin{equation}
 \int R(z_w)\,d\Pi_{S_n,\lambda}(w)
 \le\inf_{g\in\mathcal G_{C,B}}R(g)+\varepsilon.
 \label{eq:nongaussian-shallow-oracle}
\end{equation}
The risk of the ordinary raw posterior mean is no larger. For fixed
$C,H,n$, one event gives this conclusion for every $N\ge C$ and
every prefactor $\nu$ satisfying \eqref{eq:nongaussian-sample-budget}.

The learner has one prescribed architecture and one independent
coordinate prior. The comparison laws used in the proof do not
modify that prior or introduce a distribution over network sizes.
At $\nu=1$, the depth factor is $\exp(O(\sqrt H))$. The constants and
the degree of the width polynomial may depend on $d,B,\varepsilon$
and the label envelope; the theorem is not a jointly polynomial
bound in these quantities.
\end{theorem}

\paragraph{Calibration.}
Put $D_C=1+\log(2+C)$ and
$N_0=\lceil(2+C)^a(H+1)^a\rceil$.
Choose fixed positive constants $\ell_-,\ell_0,\ell_+$ sufficiently
small. If $N<N_0$, take
$\lambda=\ell_-(2+C)^{-a}(H+1)^{-a}$.
If $N\ge N_0$ and $H<aD_C$, take
$\lambda=\ell_0/[H\exp(aD_C/H)]$.
Otherwise put
\[
 T_{C,H,\nu}=1+(1-\nu)^2H+D_C^2/H,
 \qquad j=\lceil\log_2 T_{C,H,\nu}\rceil,
 \qquad\lambda=\ell_+2^{-j}.
\]
Constants can be enlarged in the order used in the proof. This
calibration depends on the declared comparison class and architecture,
not on unknown teacher coefficients or population risks.

\begin{proof}
We spell out how the Gaussian scalar construction and the preceding
new lemmas yield the claim. This also separates properties of the
comparison law from properties required of the actual prior.

\emph{A bounded comparison response.}
A fixed shallow teacher can be approximated by a bounded Gaussian
feature density:
\[
 g_0(x)=\E[\psi(G)\tanh(G\cdot(1,x))],
 \qquad G\sim N(0,I/d),\qquad \|\psi\|_\infty\le Q.
\]
Place bounded-density Gaussian bumps near its $C$ parameter vectors.
Their radii can be inverse polynomial in $C$ at the required fixed
risk accuracy. Thus $Q\le(2+C)^{a_0}$, uniformly over the class.
In particular $D=\log(e+A_dQ/\mathcal B)=O(D_C)$, where
$\mathcal B$ is a fixed auxiliary readout gain.

The constructions in Theorems~\ref{thm:quantitative-echo} and~\ref{thm:high-gain-echo} use an initial critical
variance-decreasing block, a middle gain block, and a reversed
variance block. Its scalar state is exactly
\eqref{eq:nongaussian-gaussian-state}. It supplies controls satisfying
\eqref{eq:control-product-bounds}, a bounded response
$f_{g,\nu,H}=\mathcal B r_H$, and the action estimates
\begin{align}
 \sum_{\ell=2}^H(\sqrt{W_\ell}-1)^2
 &\le A_d(1+D^2/H)\qquad(H\ge A_dD),\notag\\
 \sum_{\ell=2}^H(\sqrt{W_\ell}-1)^2
 &\le A_d[1+H\exp(A_dD/H)]\qquad(H\ge2).
 \label{eq:nongaussian-schedule-action}
\end{align}
These follow from the critical variance sequence
$q_{j+1}=\E\tanh^2(\sqrt{q_j}G)\asymp(j+1)^{-1}$:
reverse gains telescope, and a gain $e^D$ distributed over a block
of length proportional to $H$ costs $O(D^2/H)$ when $D/H$ is bounded,
or $O(H e^{A_dD/H})$ in general. The schedules for $H=2,3$ use a
fixed number of gain layers and obey the second bound.
The deterministic response has the same sign as $g_0$, has magnitude
at most $|g_0|$, and differs from it by at most
$A_d|g_0|^3/\mathcal B^2$. The scalar construction, including~\eqref{eq:continuous-cubic-response} and~\eqref{eq:response-attenuation}, is unchanged
by the present extension; its ingredients are the Gaussian state map,
the critical variance drift, and the global projection inequality
proved in Lemma~\ref{lem:nongaussian-global-transfer}.

These facts give a risk comparison without a teacher-output cap.
For $|y|\le L$ and $|g_0|\ge2L$, any attenuation of $g_0$ toward
zero cannot increase its squared loss. For $|g_0|\le2L$, choose the
fixed $\mathcal B$ so that the cubic error is sufficiently small.
For all $y$, the excess loss of attenuation is at most $y^2$.
For sub-Gaussian labels, choose fixed $L$ with small
$\E[Y^2\mathbf1_{|Y|>L}]$, and use one common empirical label-tail
and moment event. Taking the Gaussian feature approximation accuracy
inverse polynomial in $C$ controls its loss error as well. Thus the
bounded response has empirical risk at most the original teacher
risk plus the prescribed small error, on a common event.

\emph{A non-Gaussian comparison law.}
Replace only the first-layer comparison distribution by the canonical
Gaussian law. Lemma~\ref{lem:nongaussian-first-layer} charges a fixed
normalized entropy $A_{P,d}$. At subsequent layers use the adaptive
rank-one transports of Lemma~\ref{lem:nongaussian-rank-one-entropy}.
Their normalized action is bounded by
\[
 A_P\sum_\ell(\sqrt{W_\ell}-\sqrt\nu)^2
 \le 2A_P\sum_\ell(\sqrt{W_\ell}-1)^2
                   +2A_PH(1-\sqrt\nu)^2.
\]
Translate the readout mean by $\mathcal B e_H/\sqrt N$ and leave its
centered native noise unchanged. This costs at most $M\mathcal B^2$.
By Lemma~\ref{lem:nongaussian-global-transfer}, width
$N\ge(2+C)^a(H+1)^a$ gives any fixed small pointwise MSE around the
same bounded deterministic response. Consequently the normalized
comparison entropy is at most
\begin{equation}
 A[1+(1-\nu)^2H+D_C^2/H]
 \quad\hbox{if }H\ge aD_C,
 \label{eq:nongaussian-long-action}
\end{equation}
and at most $A[1+H e^{aD_C/H}]$ at smaller $H\ge2$.
The comparison is data-independent. Conditional entropy chain rules
and convexity account for its auxiliary directions.
For $H=1$, use the Gaussian first features and readout mean
$\psi(G_i)/N$. Its normalized entropy and approximation width are
polynomial in $Q$, hence polynomial in $C$; no variance schedule is
needed.

\emph{The large-width statistical bound.}
The variational free-energy inequality bounds the empirical evidence
certificate by the comparison empirical risk plus $\lambda$ times
its normalized entropy. With the stated calibrations this additional
cost is as small as desired. In the long-depth branch,
Corollary~\ref{cor:nongaussian-sobolev} requires sample size, up to
polynomial factors, at most
\[
 \exp\{A_d\sqrt{H[1+(1-\nu)^2H+D_C^2/H]}\}
 \le\exp\{A_d[\sqrt H+(1-\nu)H+D_C]\}.
\]
This has the form \eqref{eq:nongaussian-sample-budget}.
The dyadic entropy levels used in the calibration produce one sample
event for all qualifying $\nu$; their confidence costs are logarithmic.
The function classes in the evidence certificate do not depend on
learner width.

For the short-depth branch use the operator-energy certificate.
Independent sub-Gaussian rows give
\[
 \E\exp(c\omega Q_{\mathrm{op}})\le e^{AN(H+1)},\qquad
 Q_{\mathrm{op}}=N^{-1}\|W^1\|_F^2
     +\sum_{\ell=2}^H\|W^\ell\|_{\mathrm{op}}^2
     +N\|v\|_2^2.
\]
A comparison action $A_*$ therefore permits cutoff
$Q_{\mathrm{op}}=O(H+A_*)$.
The amplitude is $O(\sqrt{H+A_*})$, and AM--GM bounds the input
Lipschitz constant by
$[(H+A_*)/(H+1)]^{(H+1)/2}$.
An input-space cover has logarithmic sample budget
$O_d(H+(H+1)\log(1+A_*/H))$.
For $A_*=O(H e^{aD_C/H})$ and $H<aD_C$ this is $O_d(D_C)$,
so the required sample size is polynomial in $C$. The same reasoning
handles the $H=1$ polynomial action.

\emph{All remaining widths.}
For $C\le N<N_0$, implement the shallow teacher at the same width
and depth by weak identity relays: retain its first hidden features,
use $\eta I$ in the next layer, identity matrices subsequently, and
multiply the readout by $1/\eta$. Taking
$\eta^2=O(\varepsilon/[CB(H+1)])$ makes the accumulated tanh relay
error small. At $H=1$ use $\eta=1$ and the original teacher.
The coefficient bounds and the required parameter-box precision are
polynomial in $C,H,N$; perturbing identity layers by norm
$O(1/H)$ avoids any exponential error amplification.

The curvature upper bound implies $p(t)\ge a_Pe^{-b_Pt^2}$.
For a prior coordinate $sX$ and a target coefficient $\theta$, the
interval $|sX-\theta|\le r$ therefore has negative log probability
at most
\[
 A_P[1+\theta^2/s^2+\log_+(s/r)].
\]
After summing coordinates and multiplying by $2\nu/N$, the box law
has normalized entropy polynomial in $C,H,N$, uniformly in $\nu$;
all coordinate scales satisfy $s\le1$, and $N<N_0$ is itself
polynomial in $C,H$. Choose the small-width likelihood calibration
to pay for this comparison. Ordinary finite-parameter covers have
polynomial sample cost: parameter count is polynomial, and the
logarithm of the parameter-to-function Lipschitz constant is only
$O(H\log(C+H+N))$. Prior operator tails restore raw risk outside the
cover. A countable confidence allocation over widths and norm levels
makes this one event for the entire small-width range.

Finally, concentrate one population-risk minimizing teacher in the
compact coefficient-capped class. Its output bound $CB$ costs only a
polynomial in $C$; for sub-Gaussian labels use the corresponding
subexponential squared-loss concentration. All empirical comparison
steps and the common label-tail event then imply
\eqref{eq:nongaussian-shallow-oracle}. Jensen's inequality gives the
posterior-mean assertion. The constants may be chosen successively:
first comparison accuracies and readout gain, then polynomial width
thresholds, then the small likelihood constants, and finally the
sample exponent $c$.
\end{proof}
